% 修订正文：保留六个承担建模任务的公式；数值已展开；与独立修订稿采用同一证据。
\section{问题四的模型与求解}
\label{sec:q4r}

\subsection{问题变化与文献方法的适用边界}

第三问利用测角信息收缩源的位置可行域，并协调搜索、定位与清除。第四问增加定向发射源后，“到达一个几何上有利的测点”不再意味着能获得有效方位：无信号既可能由超出作用半径引起，也可能由机器人处于辐射背面引起。忽略这一变化会造成无效绕行；若把无信号全部解释为超距，又可能错误排除真实源。因此，需要同时回答两个问题：如何从任意朝向发现所有源，以及如何把含糊的无信号转化为可靠的位置约束。

已有文献分别提供了这两个问题的部分工具。Tokekar 与 Isler\cite{Tokekar2013}用有界测角误差刻画几何可行域，支持前文的扇形求交和保证型定位。Vander Hook 等\cite{VanderHook2014}将移动与测量时间共同纳入主动定位，提出沿主要不确定方向的垂线谨慎接近目标的策略，适合作为测点选择的文献基线。但其高斯估计和测向歧义模型不等同于本题的定向发射与同点固定误差，不能直接继承原文的性能定理。

对第四问的新因素，赵倩倩等\cite{Zhao2025}将目标发射天线的方向性纳入无人机部署，支持将源位置与朝向联合建模；其多无人机、TDOA+AOA 及连续增益模型需要改写为本题的单机器人半圆可见性。Song 等\cite{Song2012}以空间概率分布引导多个未知无线电源的搜索，支持“更新位置分布后再规划运动”的思路，但其 RSS 似然不能移用于仅有方位与接收状态的接口。本文据此保留可迁移的决策思想，重新建立本题的观测约束。

沿用第三问的区域 $\Omega$、位置域 $K$、最小包围圆中心 $c(K)$ 和半径 $\rho(K)$。以 $x$ 表示未知源位置 $g$ 的候选值，作用半径 $R\in[1000,1500]$ m，类型 $z=0,1$ 分别表示全向与定向，单位向量 $n$ 表示发射朝向。机器人在 $q$ 处的可接收状态为
\begin{equation}
 h(q;x,R,z,n)=\mathbf1\{\|q-x\|\le R\}
 \mathbf1\{z=0\ \text{或}\ n^{\mathsf T}(q-x)\ge0\}.
 \label{eq:q4r-observation}
\end{equation}
有效接收后还要满足测角误差或 5 m 近场条件。这里的发射朝向不同于测得的源方位，图\ref{fig:q4r-direction}说明了这种差异。

\begin{figure}[htbp]\centering
\includegraphics[width=\linewidth]{fig01-direction-coverage.pdf}
\caption{无信号歧义与方向鲁棒搜索原理。左图中背向和超距均产生无信号；右图中源位于站点凸包内，任意朝向下至少一站位于前向。均为机制示意，不表示官方场景真值。}
\label{fig:q4r-direction}\end{figure}

\Needspace{9\baselineskip}
\subsection{从文献基线、简洁旧版到本文策略}

为使改进具有可检验的起点，构造文献启发的谨慎贪心适配版，记为 CG。用位置域凸包的均匀面积均值 $\mu$ 和协方差近似原论文的高斯估计，令 $\sigma_{\max}$ 为最大特征值的平方根，$v$ 为对应特征向量的垂直方向。按照文献\cite{VanderHook2014}的 Algorithm 1，在 $\mu\pm av$ 中选择离当前位置较近者，其中
\begin{equation}
 a=\frac{\sigma_{\max}}{\sqrt{\sigma_{\beta_{\mathrm c}}^2-\sigma_s^2}},\qquad
 \sigma_{\beta_{\mathrm c}}=\frac{\pi}{2\Phi^{-1}(1-\beta_{\mathrm c}/2)}.
 \label{eq:q4r-cautious}
\end{equation}
测试前固定 $\beta_{\mathrm c}=0.05$，$\sigma_s^2$ 取 $\pm1^\circ$ 均匀误差的二阶矩。位置域由有界几何约束更新，清除用包围圆判断，不把协方差当作最坏误差保证。若无信号使域不变且两点均已测量，则以 $30^\circ$ 步长旋转同半径候选以避免重复测量；无可用点时转有限光学覆盖。此为公开说明适配步骤的基线，不是原论文完整 EKF 算法的复现。

CG 的局限是测点选择只考虑几何不确定性及两点之间的移动差别，没有评价“到达后能否接收到定向信号”。团队的简洁旧版 M07 在此基础上提供另一类对照：用位置域包围圆中心表示目标，以有限半径、朝向假设评估候选点的接收可能性，将失联恢复与移动代价纳入评分，并协调已知目标和搜索站点的访问次序。其不足在于，虽然无信号能影响测点评分，却没有直接收缩连续位置域；几何中心也没有反映不同候选位置对历史反馈的支持程度。

本文策略 M32 进一步维护位置、半径、类型和朝向的联合相容条件：用无信号历史排除已证伪的连续位置块，并用历史条件下的加权代表位置安排路线和测点。三者的差别见表\ref{tab:q4r-models}。为避免把缺少搜索或清除功能的弱基线作为对照，CG 与 M07、M32 共用 22 站搜索覆盖、路线组合、顺路测量、数量上界停止规则和有限清除回退；新增基线的参数在官方演练开始前冻结，未根据成绩调整。

\begin{table}[htbp]\centering\small
\caption{对照策略与逐步解决的问题}\label{tab:q4r-models}
\begin{tabularx}{\linewidth}{L{2.3cm} Y Y}\toprule
策略 & 决策依据 & 仍存在的问题或新增作用\\\midrule
文献适配 CG & 沿主要不确定方向的垂线测量，选较近侧 & 有利交会角未必可接收，可能反复失联后进入光学覆盖\\
简洁旧版 M07 & 有限假设估计可见性与行动费用，包围圆中心用于规划 & 无信号不直接裁剪连续位置域，目标代表位置较粗\\
本文 M32 & 保留 M07 框架，加入联合域证书和后验代表位置 & 将接收历史同时用于几何排除与运动规划；保证条件与启发式排序分开\\\bottomrule
\end{tabularx}\end{table}

\subsection{先保证发现：朝向鲁棒的搜索覆盖}

仅以接收圆盘覆盖目标区域不足以保证定向源被发现。对任意候选源位置 $x$，若存在站点子集 $S_x$ 满足
\begin{equation}
 x\in\operatorname{conv}(S_x),\qquad
 \max_{s\in S_x}\|s-x\|\le1000\ \mathrm m,
 \label{eq:q4r-cover}
\end{equation}
则任意朝向下均能在至少一站收到信号。证明只需将 $x$ 写成站点的凸组合：各站相对 $x$ 的有向投影的加权和为零，因此不可能所有正权重站点均严格位于背向；至少一站前向且不超出最小作用半径。

采用中心、半径 970 m 的七点等角内环及半径 1850 m 的十四点等角外环，共 22 个站点，两环均从正横轴起始。外环略超出源分布圆域，为边界源提供外侧观测位置。对覆盖半径 1800 m 圆域的 4,988 个矩形单元，逐一检查四角位于见证站点凸包内，且到各见证站点距离均小于 1000 m；最大见证距离为 999.975 m。矩形的凸性与距离函数的凸性将四角检查推广到整个单元，因而这是一份连续区域证书，而非离散源位置的抽样检验。该布局是可行配置，不主张站点数最少。

每站扫描尚未发现的频道，并保存包括无信号在内的全部反馈。访问完全部站点后，覆盖证书保证搜索完整；若提前发现 16 个不同源，则由题目数量上界提前结束搜索。随后继续处理已发现未清除源，直至全部完成。停止依据来自覆盖或数量上界，而不是连续数次无信号。

\subsection{再利用失联：位置、半径与朝向的联合裁剪}

对某频道，将与全部公开历史相容的参数记为联合域 $\Theta_t$，其位置投影为理想位置域；算法维护该投影的保守外包 $K_t$。有效方位仍用第三问的扇形求交处理，近场用 5 m 圆盘处理。新增步骤是将无信号与有效接收联合起来判断，而不是单独删除每个无信号点周围的圆盘。

设有效接收位置集合为 $\mathcal S^+$、无信号位置集合为 $\mathcal S^-$。考虑一个由圆盘 $B(c,r)$ 外包的位置块。有效接收给出整块共享的半径下界，并确定无法由超距解释的无信号点：
\begin{equation}
 L_B=\max\{1000,\max_{p\in\mathcal S^+}(\|p-c\|-r)\},\qquad
 \mathcal S^-_B=\{q\in\mathcal S^-:\|q-c\|+r<L_B\}.
 \label{eq:q4r-forced}
\end{equation}
若 $L_B>1500$，该块不可能包含源。若 $\mathcal S^-_B$ 非空，该块内任何相容源都必须为定向型，且对 $\mathcal S^-_B$ 中各点背向，同时对 $\mathcal S^+$ 中各点前向。这些限制共享同一个朝向，不能为每次观测分别任意选择解释。

对有效接收点取 $d=p-c$，对被迫背向点取 $d=c-q$。当 $\|d\|>r$ 时，块内位置变化引起的向量方向偏移至多为 $\arcsin(r/\|d\|)$；因此相容朝向必在以 $\arg d$ 为中心、半宽为 $\pi/2+\arcsin(r/\|d\|)$ 的闭圆周区间 $A(d,r)$ 内。当 $\|d\|\le r$ 时不加入该项限制。若 $\mathcal S^-_B\ne\varnothing$ 且
\begin{equation}
 \bigcap_{p\in\mathcal S^+}A(p-c,r)\ \cap\!
 \bigcap_{q\in \mathcal S^-_B}A(c-q,r)=\varnothing,
 \label{eq:q4r-reject}
\end{equation}
则整个位置块均可排除。这里求交的是必要条件的放宽，故非空时继续保留，零长度角区间也不删除。该证书将“无信号可能有多种原因”转化为“所有原因能否与同一组参数共同成立”的检验。

实现沿最小旋转外接矩形的长边二分，块尺度降至 10 m 或处理数达到 4,096 时停止；未证伪区域全部保留。矩形和圆盘留微米量级外扩，距离判据留 $10^{-5}$ m 余量，角区间再放宽 $10^{-9}$ rad，避免边界误删。图\ref{fig:q4r-domain}给出实际日志中的裁剪。所示包围圆半径由 22.42 m 降至 17.40 m，跨过 19.5 m 的直接清除阈值，说明新增信息约束能改变后续行动，而不只是改变一张位置估计图。

\begin{figure}[htbp]\centering
\includegraphics[width=\linewidth]{fig02-joint-domain.pdf}
\caption{联合约束如何促成可靠清除。左图取全部 65 次半径严格下降事件按相对降幅排序的第 33 次，橙色斜线区为排除部分；右图展示 M32 30 场官方演练的全部 113 次裁剪，65 次降低包围圆半径，半径比中位数为 0.9972。事件非独立场景，不将 113 当作性能样本量。}
\label{fig:q4r-domain}\end{figure}

\FloatBarrier
\subsection{将更新的信息转化为行动与可靠清除}

位置域负责清除保证，代表位置负责节省时间。M32 采用区域内均匀位置、均匀半径与朝向、全向和定向各半的工作先验，对各候选位置相容的半径和朝向积分赋权，以加权平均位置代替包围圆中心安排路线和生成测点。位置积分用 128 个确定性 Sobol 面积样本近似；权重不足时回退到几何代表位置。此为规划近似，既不宣称是官方生成分布，也不用于删除连续可行位置。

将剩余搜索站点与未清除目标代表位置组成节点集，从当前位置构造终点自由的开放路线。最近邻、上一轮路线更新、最远种子插入三种初始解分别经 2-opt 改善，取当前最短者。执行一个站点扫描或一个目标的完整服务后重规划。与第三问一致的是搜索与目标任务共同排序；本问用代表位置近似服务位置，不把路线长度当成全部服务时间的精确表达。

专用测点围绕代表位置及接近路径生成，包括位置域主轴的侧向点和中间点，剔除距已有测点不超过 2 m 的重复点。用至多 16 个代表位置、三个半径、24 个朝向及全向类型组成的历史相容假设，估计接收比例 $\widehat p(q)$ 和有信号后的平均包围圆半径。选择时间代理最小的点：
\begin{equation}
 \widehat J(q)=t_{\mathrm{move}}(q)+t_{\mathrm{radio}}+
 t_{\mathrm{approach}}(q)+\widehat p(q)t_{\mathrm{loc}}(q)
 +(1-\widehat p(q))t_{\mathrm{recover}}(q).
 \label{eq:q4r-decision}
\end{equation}
各项依次表示移动、一次测量、靠近目标、有信号后的继续定位及失联恢复时间代理。移动速度与测量费用取题目规定；继续定位使用剩余半径及其相对清除尺度的对数估计，恢复项使用固定基础时间和当前域尺度。冻结实现中的具体系数随可复现代码公开，属于启发式估价，不是精确的最优剩余时间。有限假设等权评分也不同于前述位置代表点的连续先验积分。

到达搜索站点时，还对仍需定位、距历史测点足够远且视角变化或接近程度足够大的已知频道顺路测量。由此一次移动可以服务搜索与定位两类需求。若 $\rho(K_t)\le19.5$ m，则移至 $c(K_t)$ 清除，20 m 光学范围覆盖整个位置域；源的发射类型不影响这一判断。每次目标服务至多进行 10 次专用无线电测量，随后用边长 25 m 的相交方格中心实施有限光学覆盖，因为方格外接圆半径仅 17.68 m。失败后保守排除已检查圆盘，在未排除候选中继续选择。有限覆盖的完整性来自几何覆盖，顺序则按单位移动及失败费用的排除面积排序。

因此，模型与有界几何计算有效时，搜索覆盖和位置域支撑完成判断；后验代表位置、路线组合与测点评分影响耗时。光学回退提供条件性完成途径，不等于所有合法场景均能在官方预算内完成。

\subsection{对照演练与同场景机制检验}

\subsubsection{三个策略各 30 次官方演练}

CG 按冻结代码新运行 30 次官方演练；M07 与 M32 各复用既有 30 场完整归档，三个样本集合共 90 个不同案例。每场先以总虚拟时间除以官方源数，再对场景取算术平均；均值区间用 29 个自由度的 Student $t$ 分布计算。场景不同，策略差异用独立样本 Welch 区间，不作同场配对。所有结果均保留，未按源数、源类型或成绩筛选。

\begin{table}[htbp]\centering\small
\caption{文献适配基线、简洁旧版与本文策略的官方演练比较}\label{tab:q4r-official}
\begin{tabularx}{\linewidth}{L{2.6cm} Z Z Z Z}\toprule
策略 & 全清场数 & 均值 & 均值 95\% 区间 & 第 90 百分位\\
 & & \multicolumn{3}{c}{时间单位：s/源}\\\midrule
文献适配 CG & 30/30 & 651.83 & [587.29, 716.36] & 825.11\\
旧版 M07 & 30/30 & 496.70 & [456.42, 536.98] & 636.08\\
本文 M32 & 30/30 & \textbf{458.37} & [416.05, 500.69] & 588.52\\\bottomrule
\end{tabularx}\end{table}

\begin{figure}[htbp]\centering
\includegraphics[width=\linewidth]{fig03-official-comparison.pdf}
\caption{全部 90 场独立官方演练及分项计费。左图展示每场结果、场均值与 95\% 均值区间；右图先对每场分项时间除以该场源数再平均。光学费用包括失败尝试，不把程序计算时间混入虚拟时间。}
\label{fig:q4r-official}\end{figure}

M32 在 30 场中清除 390 个源，均值 458.37 s/源。相对 CG，平均节省 193.46\ s/源，即 29.68\%，节省量 95\% 区间为 [117.67, 269.25]\ s/源；相对 M07，平均节省 38.33\ s/源，即 7.72\%，区间为 [-18.85, 95.51]\ s/源。本批数据支持 M32 相对文献适配版的时间改善，但相对 M07 的区间仍跨零，不能仅凭均值下降认定稳定领先。

费用与观测日志解释了改进来自哪里。CG 出现 617\ 次已发现频道的无信号测量及 928\ 次失败光学尝试；M07 对应为 393\ 和 592\ 次，M32 为 194\ 和 0 次。这里区分了正常搜索未知频道的无信号与已发现频道的失联。CG 的几何测点并非总能接收，旧版可见性评分减少部分低效行动；M32 再把历史支持程度用于位置代表点与联合域，改善剩余移动及测量。移动时间从 CG 的 509.27 s/源降至 M32 的 333.31 s/源，占两者均值差约 91\%；相对 M07，M32 的移动和无线电测量分别少 25.05 和 8.07 s/源。三组场景的平均源数为 12.77、12.87、13.00，定向源比例亦非严格匹配，因此这些频次和分项差值是机制线索，不能直接作单模块的因果归因。

\subsubsection{同一场景下，时间究竟省在哪里}

另在统一校准本地模拟器中固定 30 个新场景与同点误差，对三策略逐场比较。三者均 30/30 全清；CG、M07、M32 均值分别为 601.51、505.08、484.70 s/源。M32 比 M07 在 25 场更快、5 场更慢，平均节省 20.38 s/源，配对差值 95\% 区间为 [11.38, 29.37] s/源。这是同场诊断结果，不替代官方成绩；其场景分布与误差场仍是对官方环境的近似。

\begin{figure}[htbp]\centering
\includegraphics[width=\linewidth]{fig04-paired-paths.pdf}
\caption{三策略在同一场景的路径与清除进度。取全部 30 个本地场景按 M07 减 M32 的秒/源差值排序第 15 场，含 10 源、其中 3 个定向；圈为起点，叉为成功清除停点，非源真值。坐标等比例且范围相同；进度曲线延续至搜索完成，末端方块表示任务结束，时间计入全部付费动作。}
\label{fig:q4r-paired}\end{figure}
\FloatBarrier

图\ref{fig:q4r-paired}所示场景中，M07 与 M32 分别移动 20145.21 m 和 19346.71 m，测量 281 次和 273 次，换频 266 次和 268 次。M32 少走 798.49 m，节省 159.70 s；少测 8 次，节省 40 s；多换频 2 次，增加 2 s，合计节省 197.70 s。CG 总时间为 6112.14 s，M32 为 5552.34 s。这说明新策略的收益是较短的行动序列，而不是免费少算了若干次几何量；清除进度也显示其并非每个时刻都领先。

联合裁剪在 M32 官方日志中共排除 1,374 个区域块，113 次裁剪中 65 次降低包围圆半径，但整体半径比中位数仅为 0.9972。结合图\ref{fig:q4r-domain}，它能在部分状态直接促成清除，不能据此将全部时间提升归因于联合域。此前 M32 与使用相同后验规划的 M16 的独立官方对照，均值差仅 12.53 s/源，95\% 区间 [-41.84, 66.90] 跨零，也提示联合域的独立时间贡献仍需更强的同场消融证据。

\Needspace{18\baselineskip}
\subsection{正式测验结果}

冻结的 M32 已完成三次官方正式测验，结果见表\ref{tab:q4r-formal}，官方历史页面均显示已上传。平均定位清除时间为累计虚拟时间除以成功清除数，程序运行时间为官方接受退出与进入请求的时间戳之差，包含接口交互时间。

\begin{table}[htbp]\centering\small
\caption{M32 的三次官方正式测验结果}\label{tab:q4r-formal}
\begin{tabularx}{\linewidth}{L{5.6cm} Z Z Z}\toprule
案例编码 & 成功清除数 & 平均定位清除时间/s & 程序运行时间/s\\\midrule
9N77-TH84-ME6J-NRK4 & 10 & 557.293 & 4.325\\
QGEZ-KRWA-MNGB-6CMH & 11 & 525.395 & 4.426\\
KTZ8-EZQ5-937B-QNEV & 11 & 501.464 & 4.530\\\bottomrule
\end{tabularx}\end{table}

\FloatBarrier

正式元数据未给出实际源总数，因此仅报告已核验的成功清除数，不据此宣称官方确认全清，也不把三场正式测验并入演练置信区间。综合上述证据，M32 通过可见性约束、历史联合排除和后验代表位置衔接了第三问的几何定位与第四问的方向不确定性；其性能结论限于已测样本，工作先验、启发式恢复费用和有限计算精度仍构成适用边界。
